(1) 质量m的粒子在双势阱U(x)=lambda(x²-a²)²/4中运动，lambda,a>0。求平衡位置、稳定性、阱底小振动频率及中心势垒高度。
(2) 对能量0<E<U(0)的左阱振动，求精确转折点与用定积分表示的周期T(E)，阐明为何它可写成椭圆积分。
(3) 当E从下方接近势垒顶E_b=U(0)时，严格求周期发散的主对数系数，解释分界轨道及为何小振动近似在该极限失效。
(4) 写出x=0上方能量的跨阱轨道转折点与总能量守恒，比较E趋0和E趋E_b的物理行为，给出近势垒局部时间尺度的量纲核验。
